% \iffalse meta-comment
%
% hit-caption.dtx 是各个类共用的双语题注
%
% \fi
%
% \section{共享双语题注}
%
% \changes{v3.2a}{2026/08/07}{双语题注抽出来给各个类共用，原先每个类各写一遍}
% \changes{v3.2a}{2026/08/07}{英文题注改成 \cs{caption} 的尾随可选参数，
%   \cs{bicaption} 降为兼容写法}
% 中英文对照的题注原先要换一个命令 \cs{bicaption}，用户得记住「这张图要不要双语」
% 并挑命令。现在写法只有一个：
% \begin{latex}
% \caption{中文题}
% \caption{中文题}[English caption]
% \caption[目录用短标题]{中文题}[English caption]
% \end{latex}
% 给了尾随的 \oarg{英文题} 就排双语，不给就单语。目录用的短标题仍是 \cs{caption}
% 原本那个前置可选参数，两者位置不同，不会混。
%
% \option{style/bilingual-caption} 是一键关掉所有双语题注的开关，默认开。关掉之后
% 英文题一概不排，同一份稿子投不同用途时不必逐条删。没做成按学位给默认值：写了就出、
% 不写就不出已经够用，多一层默认反而要用户去猜自己这一档是开是关。
%
% 键自己声明，没走 \cs{\_\_hit\_style\_pair:nn}：那个在 \file{hit-key-engine.dtx} 里，
% 而本模块在 \file{hithesis.ins} 的名单里排在它前面。
%
% \cs{bicaption} 仍然可用：\pkg{ccaption} 的五参数原形与本模板之前的三、四参数写法
% 都接得住，只是各提示一次，v4 移除。
%
%    \begin{macrocode}
%<*shared-caption>
%    \end{macrocode}
%
% \changes{v3.2a}{2026/08/10}{浮动体的四个间距统一，题注间距交给 \pkg{caption} 的
%   \cs{captionsetup}，并补上内核那个前后不对称的补丁。原先三个间距互不相同、题注
%   与图表之间的空白要用户在正文里手写 \cs{vspace}，同一份文档里几张表能差出几个磅}
% 浮动体的间距照 \pkg{thuthesis} 的办法收进类里。三件事：
%
% 一、四个位置同一个值。浮动体落在正文中间、页顶页底、两个浮动体之间、整页浮动，
% 原先 \cs{intextsep} 与 \cs{textfloatsep} 是 8.5pt、\cs{floatsep} 是 12pt、
% \cs{@fpsep} 是 10pt，四种位置四个样。
%
% 二、题注与图表之间的空白由 \cs{captionsetup} 的 \texttt{skip} 给，配上
% \texttt{tableposition} 与 \texttt{figureposition}，\pkg{caption} 自己知道表题在上
% 该把空白放下面、图题在下该放上面。取 6bp，跟原先示例里手写的 \cs{vspace}\{0.5em\}
% 等价（正文 12bp 时 0.5em 是 6.02pt），所以表格那边版面不变，图那边补上原先没有的
% 那道空白。\pkg{caption} 由 \pkg{subcaption} 带进来，一直都在。
%
% 三、\hologo{LaTeX} 内核里浮动体与正文的前后距离本身不对称，光把长度设成一样没用，要打
% 三个补丁。出处见 \pkg{thuthesis} 的 \texttt{tuna/thuthesis\#614}。这段必须退出
% \pkg{expl3} 写：\cs{patchcmd} 把宏体按当前 catcode 重扫，\texttt{\string\vskip\space
% \string\intextsep} 里那个空格在 \pkg{expl3} 下会被吞，补丁就对不上（CONTRIBUTING §4.7）。
%
% 放在 \cs{AtBeginDocument} 里。\cs{AtEndOfClass} 那个时机太早：\cs{@addtocurcol}
% 那时还不是最终那一版，三个补丁全都找不到目标（实测报「没打上」），到
% \cs{begin\{document\}} 时 \cs{ifpatchable} 才为真。长度也一样，晚设才不会被覆盖。
%    \begin{macrocode}
\ExplSyntaxOff
\AtBeginDocument{%
  \setlength{\floatsep}{12bp \@plus 2bp \@minus 2bp}%
  \setlength{\textfloatsep}{12bp \@plus 2bp \@minus 2bp}%
  \setlength{\intextsep}{12bp \@plus 2bp \@minus 2bp}%
  \setlength{\@fptop}{0bp \@plus 1fil}%
  \setlength{\@fpsep}{12bp \@plus 2fil}%
  \setlength{\@fpbot}{0bp \@plus 1fil}%
  \captionsetup{skip=6bp,tableposition=top,figureposition=bottom}%
  \patchcmd{\@addtocurcol}%
    {\vskip \intextsep}%
    {\edef\hit@save@first@penalty{\the\lastpenalty}\unpenalty
     \ifnum \lastpenalty = \@M
       \unpenalty
     \else
       \penalty \hit@save@first@penalty \relax
     \fi
     \ifnum\outputpenalty <-\@Mii
       \addvspace\intextsep
       \vskip\parskip
     \else
       \addvspace\intextsep
     \fi}%
    {}{\ClassWarningNoLine{hithesis}{补丁 \string\@addtocurcol 之一没打上}}%
  \patchcmd{\@addtocurcol}%
    {\vskip\intextsep \ifnum\outputpenalty <-\@Mii \vskip -\parskip\fi}%
    {\ifnum\outputpenalty <-\@Mii
       \aftergroup\vskip\aftergroup\intextsep
       \aftergroup\nointerlineskip
     \else
       \vskip\intextsep
     \fi}%
    {}{\ClassWarningNoLine{hithesis}{补丁 \string\@addtocurcol 之二没打上}}%
  \patchcmd{\@getpen}{\@M}{\@Mi}%
    {}{\ClassWarningNoLine{hithesis}{补丁 \string\@getpen 没打上}}%
}
\ExplSyntaxOn
%    \end{macrocode}
%
% \begin{macro}{\hit@setup@caption}
% 接管 \cs{caption} 要等到 \texttt{begindocument/end}。\pkg{caption}（由
% \pkg{subcaption} 带进来）和 \pkg{ccaption} 都在 \cs{AtBeginDocument} 里重定义
% \cs{caption}，在类文件里接管会被它们覆盖掉——踩过一次，接管看似成功，
% \cs{meaning}\cs{caption} 到了正文里已经是 \pkg{caption} 的定义了。
% \texttt{begindocument/end} 排在所有 \cs{AtBeginDocument} 之后。
%
% \cs{\_\_hit\_bicaption:w} 存的是 \pkg{ccaption} 的原形，这个在类文件里存就行；
% \cs{hit@original@caption} 存的是接管前那条完整的 \cs{caption} 链。
%    \begin{macrocode}
\bool_new:N \g__hit_bilingual_caption_bool
\bool_gset_true:N \g__hit_bilingual_caption_bool
\keys_define:nn { hit / style }
  {
    bilingual-caption .bool_gset:N = \g__hit_bilingual_caption_bool ,
    bilingual-caption .default:n = { true } ,
  }
%    \end{macrocode}
%
% 排双语题注：走 \pkg{ccaption} 的原形，英文图表名取 \option{const} 族里的
% \option{figure-prefix-en} 与 \option{table-prefix-en}，用户不必每次手写
% \texttt{Fig.} 加负细空格去抵 \cs{fnum@figure} 插的那个空格。
%
% \changes{v3.2a}{2026/08/08}{双语题注的短标题接到 \pkg{ccaption} 的第二个参数上。
%   原先接的是可选参数，而 \cs{bicaption} 的可选参数是 \cs{label} 不是短标题，
%   于是 \cs{caption}\oarg{短标题} 在双语时被悄悄当成了标签}
% 短标题给 \cs{bicaption} 的\emph{第二个必选参数}，不是可选参数。ccaption 的
% \cs{bicaption}\oarg{标签}\marg{短标题}\marg{题注}\marg{英文图表名}\marg{英文题注}
% 里，可选参数是 \cs{label}：
% \begin{latex}
% \newcommand{\bicaption}[5][\@empty]{%
%   \@if@contemptyarg{#2}{\caption{#3}}{\caption[#2]{#3}}
%   \ifx\@empty#1\else \label{#1} \fi
% \end{latex}
% 接错了的后果是同一个 \cs{caption}\oarg{短标题} 两条路径含义不同：单语时进插图索引，
% 双语时变成一个标签，索引里排的还是完整题注。
%
% \changes{v3.2a}{2026/08/08}{双语题注之后补上 \cs{@currentlabel}，
%   \cs{label} 写在 \cs{caption} 后面才能拿到图表号}
% \pkg{ccaption} 之所以提供「可选参数当标签」这条路，是因为 \cs{bicaption} 把整段
% 包在 \cs{begingroup}\ldots\cs{endgroup} 里，\cs{caption} 设的 \cs{@currentlabel}
% 是局部的，出了组就没了——写在后面的 \cs{label} 只能拿到章节号。所以排完之后
% 按计数器重建一次 \cs{@currentlabel}，让 \cs{label} 与普通 \hologo{LaTeX} 一样写在
% \cs{caption} 后面。
%    \begin{macrocode}
\cs_new_protected:Npn \hit@bilingual@caption #1#2#3
  {
    \tl_if_novalue:nTF {#1}
      { \__hit_bicaption:w {} {#2} { \use:c { hit@ \@captype @prefix@en } } {#3} }
      { \__hit_bicaption:w {#1} {#2} { \use:c { hit@ \@captype @prefix@en } } {#3} }
    \protected@edef \@currentlabel
      { \use:c { p@ \@captype } \use:c { the \@captype } }
  }
\cs_new_protected:Npn \hit@single@caption #1#2
  {
    \tl_if_novalue:nTF {#1}
      { \hit@original@caption {#2} }
      { \hit@original@caption [#1] {#2} }
  }
%    \end{macrocode}
%
% \cs{caption} 的接管。给了尾随的英文题、并且开关没关死，才走双语。
%    \begin{macrocode}
\cs_new_protected:Npn \hit@setup@caption
  {
    \cs_set_eq:NN \__hit_bicaption:w \bicaption
    \AddToHook { begindocument/end }
      {
    \cs_set_eq:NN \hit@original@caption \caption
    \RenewDocumentCommand \caption { o +m o }
      {
        \IfNoValueTF {##3}
          { \hit@single@caption {##1} {##2} }
          {
            \bool_if:NTF \g__hit_bilingual_caption_bool
              { \hit@bilingual@caption {##1} {##2} {##3} }
              { \hit@single@caption {##1} {##2} }
          }
      }
%    \end{macrocode}
%
% \cs{bicaption} 的兼容层。三种写法靠「后面还跟不跟一个花括号组」区分：跟着的是
% \pkg{ccaption} 的五参数原形，抓走那一组照原形排；不跟就是本模板之前的三、四参数
% 写法。探之前先跳空格，与 \cs{@ifnextchar} 的惯例一致——旧示例里那五个参数是连着
% 写的，但用户可能换了行。题注后面正常不会紧跟花括号组。
%    \begin{macrocode}
    \RenewDocumentCommand \bicaption { o +m +m +g }
      {
        \peek_meaning_ignore_spaces:NTF \c_group_begin_token
          { \__hit_bicaption_legacy:nnnn {##1} {##2} {##3} {##4} }
          {
            \msg_warning:nn { hithesis } { bicaption-deprecated }
            \IfNoValueTF {##4}
              { \hit@bilingual@caption {##1} {##2} {##3} }
              { \hit@bilingual@caption {##1} {##3} {##4} }
          }
      }
      }
  }
\cs_new_protected:Npn \__hit_bicaption_legacy:nnnn #1#2#3#4#5
  {
    \msg_warning:nn { hithesis } { bicaption-five-args }
    \tl_if_novalue:nTF {#1}
      { \__hit_bicaption:w {#2} {#3} {#4} {#5} }
      { \__hit_bicaption:w [#1] {#2} {#3} {#4} {#5} }
  }
%    \end{macrocode}
% \end{macro}
%
%    \begin{macrocode}
%</shared-caption>
%    \end{macrocode}
%
% \Finale
